% from ega2-1-fr.tex
\oldpage[II]{82}

Pour prouver (i), on applique la définition
(\hyperref[II.4.4.2-fr]{4.4.2}) avec $\mathcal{E}=\mathcal{L}$, en notant
qu'alors $\mathbf{P}(\mathcal{E})$ s'identifie à $Y$
(\hyperref[II.4.1.4-fr]{4.1.4}).

Démontrons (iii). Il existe un $\mathcal{O}_Y$-Module quasi-cohérent
$\mathcal{E}$ et une $Y$-immersion $i:X\to\mathbf{P}(\mathcal{E})=P$ telle
que $\mathcal{L}=i^*(\mathcal{O}_P(1))$~; alors, si on pose
$\mathcal{E}'=g^*(\mathcal{E})$, $\mathcal{E}'$ est un
$\mathcal{O}_{Y'}$-Module quasi-cohérent et on a
$P'=\mathbf{P}(\mathcal{E}')=P_{(Y')}$
(\hyperref[II.4.1.3.1-fr]{4.1.3.1}), $i_{(Y')}$ est une immersion de
$X_{(Y')}$ dans $P'$ (\hyperref[I.4.3.2-fr]{I, 4.3.2}) et $\mathcal{L}'$ est
isomorphe à $(i_{(Y')})^*(\mathcal{O}_{P'}(1))$
(\hyperref[II.4.2.10-fr]{4.2.10}).

Pour prouver (iv), remarquons qu'il y a par hypothèse une $Y_i$-immersion
$r_i:X_i\to P_i=\mathbf{P}(\mathcal{E}_i)$, où $\mathcal{E}_i$ est un
$\mathcal{O}_{Y_i}$-Module quasi-cohérent, et
$\mathcal{L}_i=r_i^*(\mathcal{O}_{P_i}(1))$ ($i=1,2$)~;
$r_1\times_Sr_2$ est une $S$-immersion de $X_1\times_SX_2$ dans
$P_1\times_SP_2$ (\hyperref[I.4.3.1-fr]{I, 4.3.1}) et l'image réciproque de
$\mathcal{O}_{P_1}(1)\otimes_S\mathcal{O}_{P_2}(1)$ par cette immersion est
$\mathcal{L}_1\otimes_S\mathcal{L}_2$. D'autre part, posons
$T=Y_1\times_SY_2$, et soient $p_1,p_2$ les projections de $T$ dans $Y_1,Y_2$
respectivement. Si on pose
$P_i'=\mathbf{P}(p_i^*(\mathcal{E}_i))$ ($i=1,2$), on a par
(\hyperref[II.4.1.3.1-fr]{4.1.3.1}), $P_i'=P_i\times_{Y_i}T$ d'où
\[
  P_1'\times_TP_2'
  =(P_1\times_{Y_1}T)\times_T(P_2\times_{Y_2}T)
  =P_1\times_{Y_1}(T\times_{Y_2}P_2)
  =P_1\times_{Y_1}(Y_1\times_SP_2)
  =P_1\times_SP_2
\]
à un isomorphisme canonique près. De même, on a
$\mathcal{O}_{P_i'}(1)=\mathcal{O}_{P_i}(1)\otimes_{Y_i}\mathcal{O}_T$
(\hyperref[II.4.1.3.2-fr]{4.1.3.2}), et un calcul analogue (basé notamment
sur (\hyperref[I.9.1.9.1-fr]{I, 9.1.9.1} et
\hyperref[I.9.1.2-fr]{9.1.2})) montre que dans l'identification précédente,
$\mathcal{O}_{P_1'}(1)\otimes_T\mathcal{O}_{P_2'}(1)$ s'identifie à
$\mathcal{O}_{P_1}(1)\otimes_S\mathcal{O}_{P_2}(1)$. On peut donc considérer
$r_1\times_Sr_2$ comme une $T$-immersion de $X_1\times_SX_2$ dans
$P_1'\times_TP_2'$, l'image réciproque de
$\mathcal{O}_{P_1'}(1)\otimes_T\mathcal{O}_{P_2'}(1)$ par cette immersion
étant $\mathcal{L}_1\otimes_S\mathcal{L}_2$. On termine alors le raisonnement
comme dans (\hyperref[II.4.4.8-fr]{4.4.8}) en utilisant le morphisme de Segre.

Il reste à démontrer (ii). On peut d'abord se restreindre au cas où $Z$ est un
schéma affine, car en général il y a un recouvrement fini $(U_i)$ de $Z$ par
des ouverts affines~; si la proposition est démontrée pour
$\mathcal{K}|g^{-1}(U_i)$, $\mathcal{L}|f^{-1}(g^{-1}(U_i))$ et un entier
$n_i$, il suffira de prendre pour $n_0$ le plus grand des $n_i$ pour démontrer
la proposition pour $\mathcal{K}$ et $\mathcal{L}$
(\hyperref[II.4.4.5-fr]{4.4.5}). L'hypothèse implique que $f$ et $g$ sont des
morphismes séparés, donc $X$ et $Y$ sont des schémas quasi-compacts.

Il y a une immersion $r:X\to P=\mathbf{P}(\mathcal{E})$, où $\mathcal{E}$ est
un $\mathcal{O}_Y$-Module quasi-cohérent de type fini et
$\mathcal{L}=r^*(\mathcal{O}_P(1))$, en vertu de
(\hyperref[II.4.4.6-fr]{4.4.6}). Nous allons voir qu'il existe un
$\mathcal{O}_P$-Module $\mathcal{M}$ très ample relativement au morphisme
composé $P\xrightarrow{h}Y\xrightarrow{g}Z$, tel que $\mathcal{O}_P(1)$ soit
isomorphe à $\mathcal{M}\otimes_Y\mathcal{K}^{\otimes(-m)}$ pour un entier
$m$. Pour $n\geq m+1$,
$\mathcal{O}_P(1)\otimes_Y\mathcal{K}^{\otimes n}$ sera donc très ample pour
$Z$ en vertu de l'hypothèse et de (iv) appliqué aux morphismes $h:P\to Y$ et
$1_Y$~; comme $r$ est une immersion et que
$\mathcal{L}\otimes f^*(\mathcal{K}^{\otimes n})
=r^*(\mathcal{O}_P(1)\otimes_Y\mathcal{K}^{\otimes n})$, la conclusion
résultera de (i bis). Pour établir notre assertion sur $\mathcal{O}_P(1)$,
nous utiliserons le lemme suivant~:

\begin{lemma}[4.4.10.1]
\label{II.4.4.10.1-fr}
Soient $Z$ un schéma quasi-compact ou un préschéma dont l'espace sous-jacent
est noethérien, $g:Y\to Z$ un morphisme quasi-compact, $\mathcal{K}$ un
$\mathcal{O}_Y$-Module inversible très ample pour $g$, $\mathcal{E}$ un
$\mathcal{O}_Y$-Module quasi-cohérent de type fini. Il existe alors un entier
$m_0$ tel que, pour tout $m\geq m_0$, $\mathcal{E}$ soit isomorphe à un
quotient d'un $\mathcal{O}_Y$-Module de la forme
$g^*(\mathcal{F})\otimes\mathcal{K}^{\otimes(-m)}$, où $\mathcal{F}$ est un
$\mathcal{O}_Z$-Module quasi-cohérent de type fini (dépendant de $m$).
\end{lemma}

Ce lemme sera démontré en (\hyperref[II.4.5.10.1-fr]{4.5.10.1})~; le lecteur
pourra vérifier que (\hyperref[II.4.4.10-fr]{4.4.10}) n'est utilisé nulle part
dans le n\textsuperscript{o}~4.5.

